\subsection{JSON-Schnittstellen für Statusdaten}
\gls{json} wird als Datenformat zur Übertragung von Zustandsinformationen eingesetzt.
Es handelt sich um ein standardisiertes, menschenlesbares Format, das unabhängig von der verwendeten Programmiersprache genutzt werden kann~\cite{wehner_json_iot}.
Dadurch eignet es sich für den Austausch von Daten zwischen unterschiedlichen Systemen und Komponenten.

Aufgrund seiner einfachen und kompakten Struktur verursacht \gls{json} nur geringen Kommunikationsaufwand.
Dies ist insbesondere in eingebetteten Systemen mit begrenzten Ressourcen von Vorteil, da unnötiger Overhead vermieden wird~\cite{wehner_json_iot}.
Durch die einfache Parsbarkeit und die geringe Komplexität eignet sich das \gls{json}-Format zudem für den Einsatz in heterogenen und verteilten Systemen~\cite{wehner_json_iot}.
Damit wird eine Grundlage geschaffen, um die Wortuhr in übergeordnete Systeme einzubinden oder Statusinformationen automatisiert auszuwerten.

% Vielleicht späteres Kapitel
% In der Studienarbeit wird das \gls{json}-Format zur Bereitstellung maschinenlesbarer Statusdaten verwendet.
% Dazu gehören unter anderem Informationen zum Betriebszustand, zur Uhrzeit, zu Sensorwerten sowie zu Konfigurationsparametern.
% Die strukturierte Darstellung ermöglicht eine einfache Weiterverarbeitung der Daten durch externe Systeme.

% JSON als Datenformat
% REST-ähnliche APIs
% maschinenlesbare Daten
% Debugging
% Integration in andere Systeme möglich